\begin{textsample}{POS Dim 6 – rs – Score 12.00 – rs/rs\_em\_82.txt}  \label{ex:f6_pos_181}
5.5.1 Itinerário \textbf{Formativo} Saúde I a ) Objetivo : Compreender o conceito de saúde, identificando seus princípios determinantes na coletividade, vislumbrando possibilidades, desafios e perspectivas para melhoria da qualidade de vida ; b ) Eixos \textbf{Estruturantes} : Investigação Científica, Processos Criativos, \textbf{Mediação} e Intervenção Sociocultural e \textbf{Empreendedorismo} ; c ) Formas de Organização Metodológica : Laboratórios, Clubes, Oficinas, Observatórios, Incubadoras, Núcleos de Estudos e Núcleos de Criação Artística ; d ) Habilidades da Área focal Ciências da natureza e suas \textbf{tecnologias} ( CNT ) : a seguir estão listadas as habilidades para cada Eixo \textbf{estruturante}.
Eixo \textbf{Estruturante} : Investigação Científica nas CNT ( EMIFCNT01 ) : Investigar e analisar situações-problema e variáveis que interferem na dinâmica de fenômenos da natureza e/ou de processos tecnológicos, considerando dados e informações disponíveis em diferentes mídias, com ou sem o uso de dispositivos e \textbf{aplicativos} \textbf{digitais}. ( EMIFCNT02 ) : Levantar e \textbf{testar} hipóteses sobre variáveis que interferem na dinâmica de fenômenos da natureza e/ou de processos tecnológicos, com ou sem o uso de dispositivos e \textbf{aplicativos} \textbf{digitais}, utilizando procedimentos e linguagens adequados à investigação científica. ( EMIFCNT03 ) : Selecionar e sistematizar, com base em estudos e/ou pesquisas ( bibliográfica, exploratória, de campo, experimental etc. ) em fontes confiáveis, informações sobre a dinâmica dos fenômenos da natureza e/ou de processos tecnológicos, identificando os diversos pontos de vista e posicionando -se mediante argumentação, com o cuidado de citar as fontes dos recursos utilizados na pesquisa e buscando apresentar conclusões com o uso de diferentes mídias.
Eixo \textbf{Estruturante} : Processos Criativos nas CNT ( EMIFCNT04 ) : Reconhecer produtos e/ou processos criativos por meio de fruição, vivências e reflexão crítica sobre a dinâmica dos fenômenos naturais e/ou de processos tecnológicos, com ou sem o uso de dispositivos e \textbf{aplicativos} \textbf{digitais} ( como softwares de simulação e de realidade virtual, entre outros ). ( EMIFCNT05 ) : Selecionar e mobilizar intencionalmente recursos criativos relacionados às Ciências da Natureza para resolver \textbf{problemas} reais do ambiente e da sociedade, explorando e contrapondo diversas fontes de informação. ( EMIFCNT06 ) : Propor e \textbf{testar} soluções éticas, estéticas, criativas e inovadoras para \textbf{problemas} reais, considerando a aplicação de design de soluções e o uso de \textbf{tecnologias} \textbf{digitais}, programação e/ou pensamento computacional que apoiem a construção de protótipos, dispositivos e/ou equipamentos, com o intuito de melhorar a qualidade de vida e/ou os processos produtivos.
Eixo \textbf{Estruturante} : \textbf{Mediação} e Intervenção Sociocultural nas CNT ( EMIFCNT07 ) : Identificar e explicar questões socioculturais e \textbf{ambientais} relacionadas a fenômenos físicos, químicos e/ou biológicos. ( EMIFCNT08 ) : Selecionar e mobilizar intencionalmente conhecimentos e recursos das Ciências da Natureza para propor ações individuais e/ou coletivas de \textbf{mediação} e intervenção sobre \textbf{problemas} socioculturais e \textbf{problemas} \textbf{ambientais}. ( EMIFCNT09 ) : Propor e \textbf{testar} estratégias de \textbf{mediação} e intervenção para resolver \textbf{problemas} de natureza sociocultural e de natureza \textbf{ambiental} relacionados às Ciências da Natureza.
Eixo \textbf{Estruturante} : \textbf{Empreendedorismo} nas CNT ( EMIFCNT10 ) : Avaliar como oportunidades, conhecimentos e recursos relacionados às Ciências da Natureza podem ser utilizados na concretização de projetos pessoais ou produtivos, considerando as diversas \textbf{tecnologias} disponíveis e os impactos socioambientais. ( EMIFCNT11 ) : Selecionar e mobilizar intencionalmente conhecimentos e recursos das Ciências da Natureza para desenvolver um projeto pessoal ou um empreendimento produtivo. ( EMIFCNT12 ) : Desenvolver projetos pessoais ou produtivos, utilizando as Ciências da Natureza e suas Tecnologias para formular propostas concretas, articuladas com o projeto de vida. e ) Habilidades da Área complementar Matemática e suas \textbf{tecnologias} ( MAT ) : A seguir, estão listadas as habilidades para cada Eixo \textbf{estruturante} : Eixo \textbf{Estruturante} : Investigação Científica na MAT ( EMIFMAT01 ) : Investigar e analisar situações-problema identificando e selecionando conhecimentos matemáticos relevantes para uma dada situação, elaborando modelos para sua representação. ( EMIFMAT02 ) : Levantar e \textbf{testar} hipóteses sobre variáveis que interferem na explicação ou resolução de uma situação-problema elaborando modelos com a linguagem matemática para analisá -la e avaliar sua adequação em termos de possíveis limitações, eficiência e possibilidades de generalização. ( EMIFMAT03 ) : Selecionar e sistematizar, com base em estudos e/ou pesquisas ( bibliográfica, exploratória, de campo, experimental etc. ) em fontes confiáveis, informações sobre a contribuição da Matemática na explicação de fenômenos de natureza científica, social, profissional, cultural, de processos tecnológicos, identificando os diversos pontos de vista e posicionando -se mediante argumentação, com o cuidado de citar as fontes dos recursos utilizados na pesquisa e buscando apresentar conclusões com o uso de diferentes mídias.
Eixo \textbf{Estruturante} : Processos Criativos na MAT ( EMIFMAT04 ) : Reconhecer produtos e/ou processos criativos por meio de fruição, vivências e reflexão crítica na produção do conhecimento matemático e sua aplicação no desenvolvimento de processos tecnológicos diversos. ( EMIFMAT05 ) : Selecionar e mobilizar intencionalmente recursos criativos relacionados à Matemática para resolver \textbf{problemas} de natureza diversa, incluindo aqueles que permitam a produção de novos conhecimentos matemáticos, comunicando com precisão suas ações e reflexões relacionadas a constatações, a interpretações e a \textbf{argumentos}, bem como adequando -os às situações originais. ( EMIFMAT06 ) : Propor e \textbf{testar} soluções éticas, estéticas, criativas e inovadoras para \textbf{problemas} reais, considerando a aplicação dos conhecimentos matemáticos associados ao domínio de operações e relações matemáticas simbólicas e formais, de modo a desenvolver novas abordagens e estratégias para enfrentar novas situações.
Eixo \textbf{Estruturante} : \textbf{Mediação} e Intervenção Sociocultural na MAT ( EMIFMAT07 ) : Identificar e explicar questões socioculturais e \textbf{ambientais} aplicando conhecimentos e habilidades matemáticas para avaliar e tomar decisões em relação ao que foi observado. ( EMIFMAT08 ) : Selecionar e mobilizar intencionalmente conhecimentos e recursos matemáticos para propor ações individuais e/ou coletivas de \textbf{mediação} e intervenção sobre \textbf{problemas} socioculturais e \textbf{problemas} \textbf{ambientais}. ( EMIFMAT09 ) : Propor e \textbf{testar} estratégias de \textbf{mediação} e intervenção para resolver \textbf{problemas} de natureza sociocultural e de natureza \textbf{ambiental} relacionados à Matemática.
Eixo \textbf{Estruturante} : \textbf{Empreendedorismo} na MAT ( EMIFMAT10 ) : Avaliar como oportunidades, conhecimentos e recursos relacionados à Matemática podem ser utilizados na concretização de projetos pessoais ou produtivos, considerando as diversas \textbf{tecnologias} disponíveis e os impactos socioambientais. ( EMIFMAT11 ) : Selecionar e mobilizar intencionalmente conhecimentos e recursos da Matemática para desenvolver um projeto pessoal ou um empreendimento produtivo. ( EMIFMAT12 ) : Desenvolver projetos pessoais ou produtivos, utilizando processos e conhecimentos matemáticos para formular propostas concretas, articuladas com o projeto de vida.

% matched lemmas: ambiental, aplicativo, argumento, digital, empreendedorismo, estruturante, formativo, mediação, problema, tecnologia, testar
\end{textsample}
